\section{Classes supported on divisors}\label{sec:classes}

By \Cref{cor:Bone}, at each of the levels $114$, $44$, $52$, $102$, $33$,
$35$ and $39$ a single balanced multiset $\beta$ with $\nu([\beta])=1$
completes the standard multisets. For $33$ and $39$ the multiset $O_{M}$ of
\Cref{sec:odd} serves. For the other five we take $\beta$ to be the join of
two characters of the Fermat threefold, or of one character of the threefold
and one of the curve, and the main step is to show that the corresponding
pieces of $H^{3}$ of the threefold are supported on divisors. This follows
from a finite set of integers, a block pattern, which makes the quotient of
the threefold by the kernel of the character uniruled.

\subsection{Coniveau}
For a smooth projective variety $V$ let $N^{1}H^{k}(V,\QQ)$ be the set of
classes in $H^{k}(V,\QQ)$ that vanish on $V\setminus D$ for some divisor
$D\subseteq V$. It is a subspace, since a union of two divisors is a
divisor, and every proper closed subset of $V$ lies in a divisor.

\begin{lemma}\label{lem:N1}
Let $V$ and $W$ be smooth projective varieties.
\begin{enumerate}[label=\textup{(\alph*)}]
\item Let $f\colon V\to W$ be a surjective morphism with $\dim V=\dim W$.
Then $f^{*}N^{1}H^{k}(W,\QQ)\subseteq N^{1}H^{k}(V,\QQ)$,
$f_{*}N^{1}H^{k}(V,\QQ)\subseteq N^{1}H^{k}(W,\QQ)$, and $f_{*}f^{*}$ is
multiplication by $\deg f$.
\item There is a divisor $D\subseteq W$ such that $N^{1}H^{k}(W,\QQ)=
\iota_{*}H^{k-2}(\tilde D,\QQ)$, where $\iota\colon\tilde D\to W$ is a
resolution of $D$ followed by the inclusion.
\end{enumerate}
\end{lemma}

\begin{proof}
(a) For an open set $U\subseteq W$ the maps $f^{*}$ and $f_{*}$ commute with
restriction to $U$ and to $f^{-1}(U)$; for $f_{*}$ this is the compatibility
of proper push-forward in Borel--Moore homology, which on a smooth variety
is cohomology by Poincar\'e duality, with restriction to open sets
\cite[\S19.1]{Ful98}. If $y$ vanishes on $W\setminus D$, then $f^{*}y$
vanishes on the complement of $f^{-1}(D)$, a proper closed subset as $f$ is
surjective. If $v$ vanishes on $V\setminus E$ for a divisor $E$, then $f(E)$
has dimension less than $\dim W$, so it lies in a divisor $D$, and $f_{*}v$
vanishes on $W\setminus D$ because $f^{-1}(W\setminus D)\subseteq V\setminus
E$. By the projection formula, $f_{*}f^{*}y=f_{*}(f^{*}y\cup1)=y\cup
f_{*}1=(\deg f)\,y$.

(b) The kernels of the restriction maps $H^{k}(W,\QQ)\to H^{k}(W\setminus
D,\QQ)$ grow with $D$, and $H^{k}(W,\QQ)$ is finite-dimensional, so
$N^{1}H^{k}(W,\QQ)$ is one of these kernels. By a theorem of Deligne
\cite[Corollaire~8.2.8]{Del74}, the sequence $H^{k-2}(\tilde D,\QQ)
\xrightarrow{\iota_{*}}H^{k}(W,\QQ)\to H^{k}(W\setminus D,\QQ)$ is
exact.\footnote{The weight filtration on $H^{k}(W\setminus D)$ starts in
weight $k$, and its lowest piece is the image of $H^{k}(W)$; the kernel of
$H^{k}(W)\to H^{k}(W\setminus D)$ is the image of the cohomology with
supports $H^{k}_{D}(W)$, whose quotient of weight $k$ is the Gysin image of
$H^{k-2}(\tilde D)$, and the purity of $H^{k}(W)$ forces the kernel into
that quotient.}
\end{proof}

\begin{lemma}\label{lem:uniruledN1}
If $Y$ is a smooth projective uniruled threefold, then $N^{1}H^{3}(Y,\QQ)=
H^{3}(Y,\QQ)$.
\end{lemma}

\begin{proof}
As $Y$ is uniruled, there are a smooth projective surface $S$ and a dominant
rational map $\PP^{1}\times S\dashrightarrow Y$. Resolving its points of
indeterminacy \cite{Hir64}, we obtain a smooth projective threefold $W$, a
birational morphism $\sigma\colon W\to\PP^{1}\times S$ and a surjective
morphism $\varphi\colon W\to Y$.

First, $N^{1}H^{3}(\PP^{1}\times S)=H^{3}(\PP^{1}\times S)$. By the
K\"unneth formula, $H^{3}(\PP^{1}\times S,\QQ)$ is the sum of
$\pr^{*}H^{3}(S,\QQ)$ and $h\cup\pr^{*}H^{1}(S,\QQ)$, where $\pr$ is the
projection to $S$ and $h$ is the class of $\{p\}\times S$. The second
summand vanishes on $(\PP^{1}\setminus\{p\})\times S$. Let $C\subseteq S$ be
a smooth hyperplane section. By the Lefschetz hyperplane theorem $H_{1}(C,
\QQ)\to H_{1}(S,\QQ)$ is surjective, so by Poincar\'e duality every class of
$H^{3}(S,\QQ)$ is the Gysin image of a class of $H^{1}(C,\QQ)$, and it
vanishes on $S\setminus C$. So the first summand vanishes on $\PP^{1}\times
(S\setminus C)$.

Next, $N^{1}H^{3}(W)=H^{3}(W)$. The morphism $\sigma$ is an isomorphism over
the complement $U$ of a proper closed subset of $\PP^{1}\times S$. For
$w\in H^{3}(W,\QQ)$ the class $\sigma^{*}\sigma_{*}w-w$ vanishes on
$\sigma^{-1}(U)$, and $\sigma^{*}\sigma_{*}w$ lies in $N^{1}$ by the first
step and \Cref{lem:N1}(a); so $w\in N^{1}H^{3}(W,\QQ)$.

Finally, for $y\in H^{3}(Y,\QQ)$ we have $y=(\deg\varphi)^{-1}\varphi_{*}
\varphi^{*}y$, which lies in $N^{1}H^{3}(Y,\QQ)$ by \Cref{lem:N1}(a).
\end{proof}

\begin{lemma}\label{lem:quotientN1}
Let $X$ be a smooth projective threefold, $K$ a finite group acting on $X$,
and $\rho\colon\tilde Y\to X/K$ a resolution that is an isomorphism over the
smooth locus of $X/K$. If $N^{1}H^{3}(\tilde Y,\QQ)=H^{3}(\tilde Y,\QQ)$,
then $H^{3}(X,\QQ)^{K}\subseteq N^{1}H^{3}(X,\QQ)$.
\end{lemma}

\begin{proof}
We may assume that $K$ acts faithfully. Let $q\colon X\to X/K$ be the
quotient map, $X_{0}\subseteq X$ the dense open set of points with trivial
stabilizer, and $Y_{0}=q(X_{0})$. Then $X_{0}=q^{-1}(Y_{0})$, the variety
$Y_{0}$ is smooth, and $q\colon X_{0}\to Y_{0}$ is a finite \'etale Galois
covering with group $K$, so that $q^{*}q_{*}=\sum_{g\in K}g^{*}$ on the
cohomology of $X_{0}$. The rational map $\theta=\rho^{-1}\circ q\colon
X\dashrightarrow\tilde Y$ is a morphism on $X_{0}$. Resolving its
indeterminacy \cite{Hir64} gives a smooth projective threefold $X''$, a
birational morphism $\sigma\colon X''\to X$ that is an isomorphism over
$X_{0}$, and a surjective morphism $\tau\colon X''\to\tilde Y$ with
$\rho\tau=q\sigma$. If $\tau(p)\in\rho^{-1}(Y_{0})$, then $q\sigma(p)\in
Y_{0}$ and so $\sigma(p)\in X_{0}$. Hence $\tau^{-1}(\rho^{-1}(Y_{0}))=
\sigma^{-1}(X_{0})$, and under the isomorphisms $\sigma^{-1}(X_{0})\cong
X_{0}$ and $\rho^{-1}(Y_{0})\cong Y_{0}$ the restriction of $\tau$ is $q$.

Let $x\in H^{3}(X,\QQ)^{K}$ and $w=\sigma^{*}x$. By the compatibility with
restriction used in the proof of \Cref{lem:N1}(a), the restriction of
$\tau^{*}\tau_{*}w$ to $\sigma^{-1}(X_{0})\cong X_{0}$ is
$\sum_{g\in K}g^{*}(x|_{X_{0}})=|K|\,x|_{X_{0}}$. So $\tau^{*}\tau_{*}w-
|K|\,w$ vanishes on the complement of a proper closed subset of $X''$. Since
$\tau_{*}w\in N^{1}H^{3}(\tilde Y,\QQ)$, \Cref{lem:N1}(a) gives
$\tau^{*}\tau_{*}w\in N^{1}H^{3}(X'',\QQ)$, hence $w\in N^{1}H^{3}(X'',\QQ)$,
and $x=\sigma_{*}\sigma^{*}x=\sigma_{*}w$ lies in $N^{1}H^{3}(X,\QQ)$.
\end{proof}

\subsection{Block patterns}
Let $X=X^{3}_{m}\subseteq\PP^{4}$. Its cohomology $H^{3}(X,\QQ)$ is
primitive, and by \Cref{prop:fermatchars}(i), $H^{3}(X,\CC)$ is the direct
sum of the lines $V(a)$, $a\in\mathfrak A^{3}_{m}$, with $V(a)\subseteq
H^{3-q,q}$ for $q=\langle a\rangle-1$. As in the proof of
\Cref{prop:fermatchars}(ii), for $a\in\mathfrak A^{3}_{m}$ the sum
$\bigoplus_{t}V(ta)$ over the units $t$ is defined over $\QQ$; let
$M_{a}\subseteq H^{3}(X,\QQ)$ be the rational sub-Hodge structure with
$M_{a}\otimes\CC=\bigoplus_{t}V(ta)$.

Fix $\alpha\in\mathfrak A^{3}_{m}$, with representatives $\bar\alpha_{k}\in
\{1,\dots,m-1\}$ of its entries, and write $g\cdot\bar\alpha=\sum_{k}g_{k}
\bar\alpha_{k}$ for $g\in\ZZ^{5}$.

\begin{definition}\label{def:blockpattern}
A \emph{block pattern} for $\alpha$ consists of integers $\delta_{1},\dots,
\delta_{r}\ge1$ and $e_{kl}\ge0$ ($0\le k\le4$, $1\le l\le r$) such that
\begin{enumerate}[label=\textup{(P\arabic*)}]
\item $D=\sum_{l}e_{kl}\delta_{l}$ does not depend on $k$;
\item for every $l$, some $e_{kl}$ is positive, some $e_{kl}$ vanishes, and
the column $e_{l}=(e_{0l},\dots,e_{4l})$ has $e_{l}\cdot\bar\alpha\equiv0
\pmod m$;
\item $\sum_{l}\delta_{l}>D$.
\end{enumerate}
\end{definition}

The integers have a geometric meaning: a column $e_{l}$ with degree
$\delta_{l}$ stands for a binary form of degree $\delta_{l}$, and $D$ is the
degree of a rational curve $s\mapsto[u_{0}(s):\dots:u_{4}(s)]$ with
$u_{k}=\prod_{l}P_{l}^{e_{kl}}$ for binary forms $P_{l}$ of degree
$\delta_{l}$. Condition (P2) makes the monomials $\prod_{k}u_{k}^{\bar
\alpha_{k}}$ powers of $m$-th powers, and (P3) says that such curves move in
families large enough to sweep out the quotient below; the proof uses only
the integers.

Let $T=(\CC^{\times})^{5}/\CC^{\times}\subseteq\PP^{4}$ be the open torus.
Its character lattice is $M_{0}=\{u\in\ZZ^{5}:\ \sum_{k}u_{k}=0\}$, the
character $u$ being the function $x^{u}=\prod_{k}x_{k}^{u_{k}}$. The group
$G=\mu_{m}^{5}/\mu_{m}$ is a subgroup of $T$, and $g\in G$ multiplies $x^{u}$
by the value at $g$ of the character $u\bmod m$ of $G$. Let $K=\ker\alpha
\subseteq G$. The characters of $G$ that are trivial on $K$ are the multiples
of $\alpha$, so $T/K$ is a torus with character lattice
\[
  L_{\alpha}=\{u\in M_{0}:\ u\equiv j\bar\alpha\pmod m\ \text{for some}\
  j\in\ZZ\}.
\]
The open set $X^{\circ}=X\cap T$ is stable under $G$, and $X^{\circ}/K$ is
the hypersurface in $T/K$ defined by the $K$-invariant Laurent polynomial
\[
  f=\sum_{k=0}^{4}x^{m(e_{k}-e_{0})}=x_{0}^{-m}\sum_{k=0}^{4}x_{k}^{m},
\]
where $e_{0},\dots,e_{4}$ is the standard basis of $\ZZ^{5}$; it is a dense
open subset of $X/K$. The Newton polytope of $f$ is the simplex
$P\subseteq L_{\alpha}\otimes\RR=M_{0}\otimes\RR$ with the five vertices
$m(e_{k}-e_{0})$, which lie in $L_{\alpha}$. As the coefficients of $f$ at
these vertices do not vanish, $f$ is nondegenerate with respect to $P$
\cite[Example~5.2]{Bat24}. Following Batyrev \cite[Definition~1.2]{Bat24},
the \emph{Fine interior} $F(P)$ is the set of points $y\in P$ with
$\nu(y)-\min_{P}\nu\ge1$ for every nonzero $\nu$ in the dual lattice
$L_{\alpha}^{\vee}=\operatorname{Hom}(L_{\alpha},\ZZ)$.

\begin{lemma}\label{lem:fine}
\begin{enumerate}[label=\textup{(\alph*)}]
\item The nonzero elements of $L_{\alpha}^{\vee}$ are the functionals
$\nu_{g}(u)=(g\cdot u)/m$ with $g\in\ZZ^{5}$ not constant and
$g\cdot\bar\alpha\equiv0\pmod m$. For $y=\sum_{k}y_{k}\,m(e_{k}-e_{0})\in P$,
with $y_{k}\ge0$ and $\sum_{k}y_{k}=1$, one has
$\nu_{g}(y)-\min_{P}\nu_{g}=\sum_{k}y_{k}g_{k}-\min_{k}g_{k}$.
\item If $\alpha$ admits a block pattern, then $F(P)$ is empty.
\end{enumerate}
\end{lemma}

\begin{proof}
(a) A linear form $\nu$ on $M_{0}\otimes\RR$ that is integral on
$L_{\alpha}\supseteq mM_{0}$ has $m\nu$ integral on $M_{0}$, so $m\nu(u)=
g\cdot u$ for some $g\in\ZZ^{5}$, and $\nu\ne0$ exactly when $g$ is not
constant. The vector $u=\bar\alpha-\langle\alpha\rangle me_{0}$ lies in
$L_{\alpha}$, and $g\cdot u\equiv g\cdot\bar\alpha\pmod m$; so
$g\cdot\bar\alpha\equiv0$. Conversely, if $g\cdot\bar\alpha\equiv0$ and
$u\equiv j\bar\alpha\pmod m$, then $g\cdot u\equiv jg\cdot\bar\alpha\equiv0$,
and $\nu_{g}$ is integral on $L_{\alpha}$. Since $\nu_{g}(m(e_{k}-e_{0}))=
g_{k}-g_{0}$, we get $\nu_{g}(y)=\sum_{k}y_{k}g_{k}-g_{0}$, and the minimum
of $\nu_{g}$ on the simplex $P$ is attained at a vertex, so it is
$\min_{k}g_{k}-g_{0}$.

(b) Let $(\delta_{l},e_{kl})$ be a block pattern and $y\in F(P)$, with
barycentric coordinates $y_{k}$ as in (a). By (P2) and (a), each column
$e_{l}$ defines a nonzero $\nu_{e_{l}}\in L_{\alpha}^{\vee}$ with
$\min_{k}e_{kl}=0$, so $\sum_{k}y_{k}e_{kl}\ge1$. Multiplying by $\delta_{l}$,
adding over $l$ and using (P1),
\[
  \sum_{l}\delta_{l}\le\sum_{l}\delta_{l}\sum_{k}y_{k}e_{kl}=\sum_{k}y_{k}
  \sum_{l}e_{kl}\delta_{l}=D\sum_{k}y_{k}=D,
\]
which contradicts (P3).
\end{proof}

\begin{theorem}\label{thm:uniruled}
Let $\alpha\in\mathfrak A^{3}_{m}$ admit a block pattern, and let
$K=\ker\alpha\subseteq G$. Then every smooth projective variety birational to
$X^{3}_{m}/K$ is uniruled, and $M_{\alpha}\subseteq N^{1}H^{3}(X^{3}_{m},
\QQ)$.
\end{theorem}

\begin{proof}
Let $\tilde Y$ be a smooth projective variety birational to $X/K$. By
\Cref{lem:fine}(b) and Batyrev's theorem \cite[Corollary~5.6]{Bat24}, the
nondegenerate hypersurface $X^{\circ}/K$ has negative Kodaira dimension;
since $X^{\circ}/K$ is birational to $\tilde Y$, so has $\tilde Y$. A smooth
projective threefold of negative Kodaira dimension is uniruled: by Mori's
theorem \cite{Mor88} it is birational to a projective threefold $Y'$ with
terminal singularities that is either minimal, with $K_{Y'}$ nef, or the
total space of a Mori fibre space. In the first case $Y'$ would have
nonnegative Kodaira dimension by Miyaoka's theorem \cite{Miy88}. In the
second, a general fibre $F$ has $-K_{F}$ ample, so $K_{F}$ is negative on
the curves cut out on $F$ by general ample divisors, and $F$ is uniruled by
the criterion of Miyaoka and Mori \cite{MM86}; hence $Y'$ is uniruled, and so
is $\tilde Y$.

Choose $\tilde Y$ to be a resolution of $X/K$ that is an isomorphism over
the smooth locus \cite{Hir64}. By \Cref{lem:uniruledN1},
$N^{1}H^{3}(\tilde Y,\QQ)=H^{3}(\tilde Y,\QQ)$, and by
\Cref{lem:quotientN1}, $H^{3}(X,\QQ)^{K}\subseteq N^{1}H^{3}(X,\QQ)$. For a
unit $t$ the character $t\alpha$ is trivial on $K$, so $K$ acts trivially on
$V(t\alpha)$ and on $M_{\alpha}$, and $M_{\alpha}\subseteq H^{3}(X,\QQ)^{K}$.
\end{proof}

\subsection{Joins}
The joins below use the inductive structure of \Cref{sec:fermat}: for
$r,s\ge1$ the map $f\colon H^{r+s-2}(X^{r-1}_{m}\times X^{s-1}_{m})\to
H^{r+s}(X^{r+s}_{m})$ carries $V(a')\otimes V(b')$ onto $V(a'*b')$ and
algebraic classes to algebraic classes.

\begin{proposition}\label{prop:joindivisor}
Let $\alpha,\alpha'\in\mathfrak A^{3}_{m}$ with $M_{\alpha}$ and
$M_{\alpha'}$ contained in $N^{1}H^{3}(X^{3}_{m},\QQ)$.
\begin{enumerate}[label=\textup{(\alph*)}]
\item If $\langle t\alpha\rangle+\langle t\alpha'\rangle=5$ for every unit
$t$, then the multiset of the ten entries of $\alpha$ and $\alpha'$ is
balanced and algebraic.
\item If $b\in\mathfrak A^{1}_{m}$ and $\langle t\alpha\rangle+\langle
tb\rangle=4$ for every unit $t$, then the multiset of the eight entries of
$\alpha$ and $b$ is balanced and algebraic.
\end{enumerate}
\end{proposition}

\begin{proof}
(a) Let $X=X^{3}_{m}$. The multiset is balanced, because for every unit $t$
the representatives of the entries of $t\alpha$ and $t\alpha'$ add up to
$5m$. For every $t$, one of $V(t\alpha)$, $V(t\alpha')$ lies in $H^{2,1}$
and the other in $H^{1,2}$, as $\langle ta\rangle\in\{1,\dots,4\}$ and
$\langle-ta\rangle=5-\langle ta\rangle$ force both values into $\{2,3\}$. So
$L=\bigoplus_{t}V(t\alpha)\otimes V(t\alpha')\subseteq H^{6}(X\times X,\CC)$
has type $(3,3)$. It is stable under $\operatorname{Gal}(\QQ(\mu_{m})/\QQ)$,
so $L=L_{\QQ}\otimes\CC$ for a space $L_{\QQ}\subseteq M_{\alpha}\otimes
M_{\alpha'}$ of Hodge classes.

By \Cref{lem:N1}(b) there are smooth projective surfaces $\tilde D,\tilde D'$
with morphisms $\iota\colon\tilde D\to X$, $\iota'\colon\tilde D'\to X$ such
that $M_{\alpha}\subseteq\iota_{*}H^{1}(\tilde D,\QQ)$ and $M_{\alpha'}
\subseteq\iota'_{*}H^{1}(\tilde D',\QQ)$. The map
\[
  \gamma=(\iota\times\iota')_{*}\colon H^{1}(\tilde D,\QQ)\otimes
  H^{1}(\tilde D',\QQ)(-2)\to H^{3}(X,\QQ)\otimes H^{3}(X,\QQ)
\]
is a morphism of polarizable Hodge structures of weight $6$, given by
$\pm\iota_{*}\otimes\iota'_{*}$ in the K\"unneth decomposition, and its image
contains $M_{\alpha}\otimes M_{\alpha'}$. Applying \Cref{lem:lift} to the
restriction of $\gamma$ to $\gamma^{-1}(M_{\alpha}\otimes M_{\alpha'})$,
every $\upsilon\in L_{\QQ}$ is $\gamma(\upsilon')$ for a Hodge class
$\upsilon'$ in $H^{1}(\tilde D)\otimes H^{1}(\tilde D')(-2)$, that is, a
Hodge class of degree $2$ on $\tilde D\times\tilde D'$. By \Cref{thm:lef11}
it is algebraic, and so is $\upsilon=(\iota\times\iota')_{*}\upsilon'$.
Hence $V(\alpha)\otimes V(\alpha')\subseteq\Alg^{3}(X\times X)\otimes\CC$,
and by the inductive structure with $r=s=4$ the line $V(\alpha*\alpha')=
f\bigl(V(\alpha)\otimes V(\alpha')\bigr)$ is algebraic.

(b) Let $X'=X^{1}_{m}$. The multiset is balanced because the entries of
$t\alpha$ and $tb$ add up to $4m$. As $\langle tb\rangle\in\{1,2\}$, we have
$\langle t\alpha\rangle\in\{2,3\}$, and $V(t\alpha)\subseteq H^{2,1}(X)$
exactly when $V(tb)\subseteq H^{0,1}(X')$; so $\bigoplus_{t}V(t\alpha)
\otimes V(tb)\subseteq H^{4}(X\times X',\CC)$ has type $(2,2)$, and as in (a)
it is $L_{\QQ}\otimes\CC$ for a space $L_{\QQ}\subseteq M_{\alpha}\otimes
H^{1}(X',\QQ)$ of Hodge classes. With $\tilde D$ and $\iota$ as in (a), the
map $(\iota\times\mathrm{id})_{*}\colon H^{1}(\tilde D,\QQ)\otimes
H^{1}(X',\QQ)(-1)\to H^{3}(X,\QQ)\otimes H^{1}(X',\QQ)$ has image containing
$M_{\alpha}\otimes H^{1}(X',\QQ)$, and \Cref{lem:lift} lifts every
$\upsilon\in L_{\QQ}$ to a Hodge class of degree $2$ on the threefold
$\tilde D\times X'$. It is algebraic by \Cref{thm:lef11}, so $\upsilon$ is
algebraic, and the inductive structure with $r=4$ and $s=2$ shows that
$V(\alpha*b)$ is algebraic.
\end{proof}

\subsection{Five pairs of characters}
For a character $a$ of $X^{3}_{m}$ or $X^{1}_{m}$ we have $\langle-ta\rangle
=5-\langle ta\rangle$, respectively $3-\langle ta\rangle$, so it suffices to
list the values at the units $t<m/2$. Each value is the sum of the
representatives in $\{1,\dots,m-1\}$ of the entries of $ta$, divided by $m$.

\begin{lemma}\label{lem:charpairs}
In each line of the following list, the characters named satisfy the
hypothesis of \Cref{prop:joindivisor}(a) or~(b), and the multiset $\beta$ of
their entries is balanced, contains no pair, and has $\nu([\beta])=1$ for the
homomorphism $\nu$ of \Cref{cor:Bone}.
\begin{enumerate}[label=\textup{(\roman*)}]
\item $m=114$: $\alpha_{1}=(1,25,43,57,102)$, $\alpha_{2}=(39,63,68,80,92)$,
$\beta=\{1,\allowbreak 25,\allowbreak 39,\allowbreak 43,\allowbreak 57,\allowbreak 63,\allowbreak 68,\allowbreak 80,\allowbreak 92,\allowbreak 102\}$, $\nu=\nu_{19}$.
\item $m=44$: $\alpha=(1,5,8,33,41)$, $b=(9,37,42)$,
$\beta=\{1,\allowbreak 5,\allowbreak 8,\allowbreak 9,\allowbreak 33,\allowbreak 37,\allowbreak 41,\allowbreak 42\}$, $\nu=\nu_{11}$.
\item $m=52$: $\alpha=(1,5,16,37,45)$, $\alpha'=(29,18,42,50,17)$,
$\beta=\{1,\allowbreak 5,\allowbreak 16,\allowbreak 17,\allowbreak 18,\allowbreak 29,\allowbreak 37,\allowbreak 42,\allowbreak 45,\allowbreak 50\}$, $\nu=\nu_{13}$.
\item $m=102$: $\alpha=(1,7,51,60,85)$, $\alpha'=(91,34,26,98,57)$,
$\beta=\{1,\allowbreak 7,\allowbreak 26,\allowbreak 34,\allowbreak 51,\allowbreak 57,\allowbreak 60,\allowbreak 85,\allowbreak 91,\allowbreak 98\}$, $\nu=\nu_{17}$.
\item $m=35$: $\alpha=(1,2,16,21,30)$, $b=(17,22,31)$,
$\beta=\{1,\allowbreak 2,\allowbreak 16,\allowbreak 17,\allowbreak 21,\allowbreak 22,\allowbreak 30,\allowbreak 31\}$, $\nu=\nu_{R}$.
\end{enumerate}
\end{lemma}

\begin{proof}
The entries are nonzero and the entry sums are multiples of $m$: $228$ and
$342$ in (i), $88$ and $88$ in (ii), $104$ and $156$ in (iii), $204$ and
$306$ in (iv), $70$ and $70$ in (v). The values at the units $t<m/2$ are:

\smallskip
\noindent(i) $m=114$:
\begin{center}
\small
\setlength{\tabcolsep}{3.2pt}
\begin{tabular}{l*{18}{c}}
\toprule
$t$ & 1 & 5 & 7 & 11 & 13 & 17 & 23 & 25 & 29 & 31 & 35 & 37 & 41 & 43 & 47 & 49 & 53 & 55\\
\midrule
$\langle t\alpha_{1}\rangle$ & 2&2&2&2&3&2&2&2&3&3&2&2&3&2&2&3&3&2\\
$\langle t\alpha_{2}\rangle$ & 3&3&3&3&2&3&3&3&2&2&3&3&2&3&3&2&2&3\\
\bottomrule
\end{tabular}
\end{center}
(ii) $m=44$:
\begin{center}
\small
\begin{tabular}{l*{10}{c}}
\toprule
$t$ & 1 & 3 & 5 & 7 & 9 & 13 & 15 & 17 & 19 & 21\\
\midrule
$\langle t\alpha\rangle$ & 2&2&3&2&2&2&3&3&2&3\\
$\langle tb\rangle$ & 2&2&1&2&2&2&1&1&2&1\\
\bottomrule
\end{tabular}
\end{center}
(iii) $m=52$:
\begin{center}
\small
\begin{tabular}{l*{12}{c}}
\toprule
$t$ & 1 & 3 & 5 & 7 & 9 & 11 & 15 & 17 & 19 & 21 & 23 & 25\\
\midrule
$\langle t\alpha\rangle$ & 2&2&2&2&3&2&3&2&3&2&2&3\\
$\langle t\alpha'\rangle$ & 3&3&3&3&2&3&2&3&2&3&3&2\\
\bottomrule
\end{tabular}
\end{center}
(iv) $m=102$:
\begin{center}
\small
\setlength{\tabcolsep}{4pt}
\begin{tabular}{l*{16}{c}}
\toprule
$t$ & 1 & 5 & 7 & 11 & 13 & 19 & 23 & 25 & 29 & 31 & 35 & 37 & 41 & 43 & 47 &
49\\
\midrule
$\langle t\alpha\rangle$ & 2&2&2&2&3&2&2&3&2&2&2&3&2&3&2&3\\
$\langle t\alpha'\rangle$ & 3&3&3&3&2&3&3&2&3&3&3&2&3&2&3&2\\
\bottomrule
\end{tabular}
\end{center}
(v) $m=35$:
\begin{center}
\small
\begin{tabular}{l*{12}{c}}
\toprule
$t$ & 1 & 2 & 3 & 4 & 6 & 8 & 9 & 11 & 12 & 13 & 16 & 17\\
\midrule
$\langle t\alpha\rangle$ & 2&2&2&2&2&3&2&2&2&3&3&3\\
$\langle tb\rangle$ & 2&2&2&2&2&1&2&2&2&1&1&1\\
\bottomrule
\end{tabular}
\end{center}
For instance, in (v), $8\alpha=(8,16,23,28,30)$ and $8b=(31,1,3)$ have
entry sums $105$ and $35$.

In each line the representatives of the entries of $t\beta$ add up to
$m(\langle t\alpha\rangle+\langle t\alpha'\rangle)$ or $m(\langle t\alpha
\rangle+\langle tb\rangle)$, a constant, so $\beta$ is balanced; and no two
entries of $\beta$ add up to $m$. The elements of order $19$ in $\ZZ/114$ are
the nonzero multiples of $6$, and $102$ is the only entry of $\beta$ of this
kind in (i); the elements of order $11$ in $\ZZ/44$ and of order $13$ in
$\ZZ/52$ are the nonzero multiples of $4$, and $8$, respectively $16$, is
the only such entry in (ii), (iii); the elements of order $17$ in $\ZZ/102$
are the nonzero multiples of $6$, and $60$ is the only such entry in (iv).
In (v) the elements of $\ZZ/35$ with square $1$ are $1$, $6$, $29$ and
$34$, and only $1$ is an entry of $\beta$.
\end{proof}

\begin{corollary}\label{cor:patterns}
Each of the characters in the following table admits the block pattern
listed, a column $e_{l}$ with degree $\delta_{l}$ being written
$e_{l}{:}\delta_{l}$:
\begin{center}
\small
\setlength{\tabcolsep}{4pt}
\begin{tabular}{@{}llc>{\raggedright\arraybackslash}p{0.52\textwidth}@{}}
\toprule
$m$ & $\alpha$ & $D$ & block pattern\\
\midrule
$114$ & $13\alpha_{2}=(51,21,86,14,56)$ & $3$ & $(0,0,1,2,0){:}1$,
$(0,0,2,0,1){:}1$, $(1,3,0,0,0){:}1$, $(2,0,0,1,2){:}1$\\
$114$ & $\alpha_{1}=(1,25,43,57,102)$ & $24$ & $(3,1,2,0,0){:}7$,
$(0,1,5,0,1){:}2$, $(0,0,0,2,0){:}12$, $(1,5,0,0,1){:}3$,
$(0,0,0,0,19){:}1$\\
$44$ & $(1,5,8,33,41)$ & $11$ & $(3,0,0,0,1){:}1$, $(0,0,11,0,0){:}1$,
$(1,1,0,0,2){:}5$, $(1,2,0,1,0){:}3$, $(0,0,0,4,0){:}2$\\
$52$ & $(1,5,16,37,45)$ & $14$ & $(0,0,0,3,1){:}4$, $(0,1,1,0,3){:}1$,
$(2,1,0,0,1){:}7$, $(0,0,13,0,0){:}1$, $(0,3,0,1,0){:}2$\\
$52$ & $9\alpha'=(1,6,14,34,49)$ & $5$ & $(0,0,5,1,0){:}1$,
$(3,0,0,0,1){:}1$, $(2,0,0,3,0){:}1$, $(0,3,0,1,0){:}1$,
$(0,1,0,0,2){:}2$\\
$102$ & $37\alpha'=(1,34,44,56,69)$ & $6$ & $(0,0,0,3,2){:}1$,
$(0,0,3,0,4){:}1$, $(0,3,0,0,0){:}2$, $(2,0,1,1,0){:}3$\\
$102$ & $5\alpha=(5,35,51,96,17)$ & $18$ & $(0,0,0,0,6){:}2$,
$(0,0,0,17,0){:}1$, $(0,0,2,0,0){:}9$, $(0,6,0,1,0){:}1$,
$(3,2,0,0,1){:}6$\\
$35$ & $(1,2,16,21,30)$ & $50$ & $(0,0,0,5,0){:}10$, $(1,1,2,0,0){:}25$,
$(1,2,0,0,1){:}10$, $(3,1,0,0,1){:}5$, $(0,0,0,0,35){:}1$\\
\bottomrule
\end{tabular}
\end{center}
In each row the degrees add up to $D+1$. Hence $M_{\alpha}\subseteq N^{1}H^{3}
(X^{3}_{m},\QQ)$ for each of these characters.
\end{corollary}

\begin{proof}
Condition (P1) says that each of the five coordinates of $\sum_{l}\delta_{l}
e_{l}$ equals $D$; for instance, in the second row these coordinates are
$21+3$, $7+2+15$, $14+10$, $24$ and $2+3+19$, and in the last row they are
$25+10+15$, $25+20+5$, $50$, $50$ and $10+5+35$. Condition (P2) is the
divisibility of the weights $e_{l}\cdot\bar\alpha$ by $m$. Row by row, the
weights are
\[
\begin{array}{@{}l@{\qquad}l@{}}
  114,\ 228,\ 114,\ 228; & 114,\ 342,\ 114,\ 228,\ 1938;\\
  44,\ 88,\ 88,\ 44,\ 132; & 156,\ 156,\ 52,\ 208,\ 52;\\
  104,\ 52,\ 104,\ 52,\ 104; & 306,\ 408,\ 102,\ 102;\\
  102,\ 1632,\ 102,\ 306,\ 102; & 105,\ 35,\ 35,\ 35,\ 1050;
\end{array}
\]
for instance $(0,0,1,2,0)\cdot(51,21,86,14,56)=86+28$. The degrees add up to
$4$, $25$, $12$, $15$, $6$, $7$, $19$ and $51$, which is (P3).
\Cref{thm:uniruled} gives the last assertion.
\end{proof}

\begin{theorem}\label{thm:classes}
The levels $33$, $35$, $39$, $44$, $52$, $102$ and $114$ are complete, and so
are $51$ and $57$.
\end{theorem}

\begin{proof}
For $M\in\{33,39\}$, $O_{M}$ is balanced and algebraic by
\Cref{lem:exceptional}. The elements of $\ZZ/33$ with square $1$ are $1$,
$10$, $23$, $32$, and those of $\ZZ/39$ are $1$, $14$, $25$, $38$; of these,
$O_{33}=\{1,4,16,22,25,31\}$ and $O_{39}=\{1,7,16,22,34,37\}$ contain only
$1$. So $\nu_{R}([O_{M}])=1$, and by \Cref{cor:Bone}, $B_{M}=S_{M}+
\ZZ[O_{M}]$.

For the other five levels let $\beta$ be the multiset of \Cref{lem:charpairs}.
Since $M_{ta}=M_{a}$ for every unit $t$, \Cref{cor:patterns} shows that
$M_{\alpha}$, and $M_{\alpha'}$ in (i), (iii), (iv), lie in
$N^{1}H^{3}(X^{3}_{m},\QQ)$: in (i), $M_{\alpha_{2}}=M_{13\alpha_{2}}$; in
(iii), $M_{\alpha'}=M_{9\alpha'}$; in (iv), $M_{\alpha}=M_{5\alpha}$ and
$M_{\alpha'}=M_{37\alpha'}$. By \Cref{prop:joindivisor}, $\beta$ is
algebraic, and by \Cref{cor:Bone}, $B_{m}=S_{m}+\ZZ[\beta]$.

In all seven cases \Cref{lem:transfer} shows that the level is complete.
The levels $57$ and $51$ divide $114$ and $102$.
\end{proof}

\begin{corollary}\label{cor:degrees}
The following levels $M$ are complete, and the Hodge conjecture holds for
$X^{n}_{M}$ for every $n$:
\begin{enumerate}[label=\textup{(\roman*)}]
\item $M=2^{a}3^{b}5^{c}7^{d}$ with $420\nmid M$;
\item $M=2^{a}3^{b}11^{c}$ and $M=2^{a}3^{b}13^{c}$;
\item $M=2^{a}3^{b}17^{c}$ and $M=2^{a}3^{b}19^{c}$ with $a\le1$.
\end{enumerate}
\end{corollary}

\begin{proof}
For $M\le2$, $X^{n}_{M}$ is a linear space or a quadric, and every
balanced multiset at level $2$ is a sum of pairs $\{1,1\}$; let $M\ge3$. In
(i), the elements of $\mathsf P$ dividing $M$ are among $4,3,5,7$, so
$\mathcal Q(M)$ is contained in $\{12,15,20,21,28,35,420\}$, and $420\nmid M$
removes the last; in (ii) it is contained in $\{12,33,44\}$ or $\{12,39,52\}$;
in (iii), where $4\nmid M$, in $\{51\}$ or $\{57\}$. These levels are
complete by \Cref{prop:aokidegrees} and \Cref{thm:classes}, and
\Cref{prop:complete} applies.
\end{proof}

Part (i) for the degrees divisible by $35$ is due to Peterson \cite{Pet26},
who proves the case $M=35$ with the character $\alpha$ of
\Cref{lem:charpairs}(v) by a family of curves on $X^{3}_{35}$ and a residue
computation; the block pattern replaces that computation. Parts (ii) and
(iii) are \cite[Theorem~A]{BBe84hc}, where the coniveau step is also proved
by families of curves.
